% Options for packages loaded elsewhere
% Options for packages loaded elsewhere
\PassOptionsToPackage{unicode}{hyperref}
\PassOptionsToPackage{hyphens}{url}
\PassOptionsToPackage{dvipsnames,svgnames,x11names}{xcolor}
%
\documentclass[
  letterpaper,
  DIV=11,
  numbers=noendperiod]{scrartcl}
\usepackage{xcolor}
\usepackage{amsmath,amssymb}
\setcounter{secnumdepth}{5}
\usepackage{iftex}
\ifPDFTeX
  \usepackage[T1]{fontenc}
  \usepackage[utf8]{inputenc}
  \usepackage{textcomp} % provide euro and other symbols
\else % if luatex or xetex
  \usepackage{unicode-math} % this also loads fontspec
  \defaultfontfeatures{Scale=MatchLowercase}
  \defaultfontfeatures[\rmfamily]{Ligatures=TeX,Scale=1}
\fi
\usepackage{lmodern}
\ifPDFTeX\else
  % xetex/luatex font selection
\fi
% Use upquote if available, for straight quotes in verbatim environments
\IfFileExists{upquote.sty}{\usepackage{upquote}}{}
\IfFileExists{microtype.sty}{% use microtype if available
  \usepackage[]{microtype}
  \UseMicrotypeSet[protrusion]{basicmath} % disable protrusion for tt fonts
}{}
\makeatletter
\@ifundefined{KOMAClassName}{% if non-KOMA class
  \IfFileExists{parskip.sty}{%
    \usepackage{parskip}
  }{% else
    \setlength{\parindent}{0pt}
    \setlength{\parskip}{6pt plus 2pt minus 1pt}}
}{% if KOMA class
  \KOMAoptions{parskip=half}}
\makeatother
% Make \paragraph and \subparagraph free-standing
\makeatletter
\ifx\paragraph\undefined\else
  \let\oldparagraph\paragraph
  \renewcommand{\paragraph}{
    \@ifstar
      \xxxParagraphStar
      \xxxParagraphNoStar
  }
  \newcommand{\xxxParagraphStar}[1]{\oldparagraph*{#1}\mbox{}}
  \newcommand{\xxxParagraphNoStar}[1]{\oldparagraph{#1}\mbox{}}
\fi
\ifx\subparagraph\undefined\else
  \let\oldsubparagraph\subparagraph
  \renewcommand{\subparagraph}{
    \@ifstar
      \xxxSubParagraphStar
      \xxxSubParagraphNoStar
  }
  \newcommand{\xxxSubParagraphStar}[1]{\oldsubparagraph*{#1}\mbox{}}
  \newcommand{\xxxSubParagraphNoStar}[1]{\oldsubparagraph{#1}\mbox{}}
\fi
\makeatother


\usepackage{longtable,booktabs,array}
\newcounter{none} % for unnumbered tables
\usepackage{calc} % for calculating minipage widths
% Correct order of tables after \paragraph or \subparagraph
\usepackage{etoolbox}
\makeatletter
\patchcmd\longtable{\par}{\if@noskipsec\mbox{}\fi\par}{}{}
\makeatother
% Allow footnotes in longtable head/foot
\IfFileExists{footnotehyper.sty}{\usepackage{footnotehyper}}{\usepackage{footnote}}
\makesavenoteenv{longtable}
\usepackage{graphicx}
\makeatletter
\newsavebox\pandoc@box
\newcommand*\pandocbounded[1]{% scales image to fit in text height/width
  \sbox\pandoc@box{#1}%
  \Gscale@div\@tempa{\textheight}{\dimexpr\ht\pandoc@box+\dp\pandoc@box\relax}%
  \Gscale@div\@tempb{\linewidth}{\wd\pandoc@box}%
  \ifdim\@tempb\p@<\@tempa\p@\let\@tempa\@tempb\fi% select the smaller of both
  \ifdim\@tempa\p@<\p@\scalebox{\@tempa}{\usebox\pandoc@box}%
  \else\usebox{\pandoc@box}%
  \fi%
}
% Set default figure placement to htbp
\def\fps@figure{htbp}
\makeatother





\setlength{\emergencystretch}{3em} % prevent overfull lines

\providecommand{\tightlist}{%
  \setlength{\itemsep}{0pt}\setlength{\parskip}{0pt}}



 


\KOMAoption{captions}{tableheading}
\makeatletter
\@ifpackageloaded{caption}{}{\usepackage{caption}}
\AtBeginDocument{%
\ifdefined\contentsname
  \renewcommand*\contentsname{Table of contents}
\else
  \newcommand\contentsname{Table of contents}
\fi
\ifdefined\listfigurename
  \renewcommand*\listfigurename{List of Figures}
\else
  \newcommand\listfigurename{List of Figures}
\fi
\ifdefined\listtablename
  \renewcommand*\listtablename{List of Tables}
\else
  \newcommand\listtablename{List of Tables}
\fi
\ifdefined\figurename
  \renewcommand*\figurename{Figure}
\else
  \newcommand\figurename{Figure}
\fi
\ifdefined\tablename
  \renewcommand*\tablename{Table}
\else
  \newcommand\tablename{Table}
\fi
}
\@ifpackageloaded{float}{}{\usepackage{float}}
\floatstyle{ruled}
\@ifundefined{c@chapter}{\newfloat{codelisting}{h}{lop}}{\newfloat{codelisting}{h}{lop}[chapter]}
\floatname{codelisting}{Listing}
\newcommand*\listoflistings{\listof{codelisting}{List of Listings}}
\makeatother
\makeatletter
\makeatother
\makeatletter
\@ifpackageloaded{caption}{}{\usepackage{caption}}
\@ifpackageloaded{subcaption}{}{\usepackage{subcaption}}
\makeatother
\usepackage{bookmark}
\IfFileExists{xurl.sty}{\usepackage{xurl}}{} % add URL line breaks if available
\urlstyle{same}
\hypersetup{
  colorlinks=true,
  linkcolor={blue},
  filecolor={Maroon},
  citecolor={Blue},
  urlcolor={Blue},
  pdfcreator={LaTeX via pandoc}}


\author{}
\date{}
\begin{document}


\includegraphics[width=0.6in,height=0.60333in]{media/image1.jpeg}\textbf{UNIVERSIDAD
NACIONAL DE LOJA}

\textbf{\emph{Carrera} \emph{de} \emph{Economía}}

\textbf{Desarrollo Económico}

\textbf{\emph{Séptimo A}}

\textbf{Unidad 1}

\textbf{Septiembre 2025/Febrero 2026}

Nombre: Fecha:

\textbf{Taller 2: INFORME DEL TALLER TEATRAL: TEORÍAS DEL DESARROLLO
ECONÓMICO}

\textbf{Tema:} (Nombre de la teoría asignada)

\subsection{1. Definición}\label{definiciuxf3n}

(Explicar en un párrafo breve y claro en qué consiste la teoría. Se
sugiere entre 5 a 6 líneas).

\subsection{2. Principales Características
Explicadas}\label{principales-caracteruxedsticas-explicadas}

\begin{itemize}
\item
  Característica 1: (Descripción)
\item
  Característica 2: (Descripción)
\item
  Característica 3: (Descripción)
\end{itemize}

(Pueden incluir de 3 a 5 características más relevantes).

\subsection{3. Políticas sugeridas}\label{poluxedticas-sugeridas}

\begin{itemize}
\item
  Política 1: (Descripción)
\item
  Política 2: (Descripción)
\item
  Política 3: (Descripción)
\end{itemize}

(Se sugiere indicar entre 2 y 4 políticas económicas que la teoría
respalde o proponga).

\subsection{4. Guion de la Obra de
Teatro}\label{guion-de-la-obra-de-teatro}

\textbf{Título de la obra:} (Proponer un título creativo)

\textbf{Escenario:} (Describir brevemente dónde y cuándo se desarrolla
la historia)

\textbf{Personajes:}

\begin{itemize}
\item
  (Nombre del personaje y breve descripción de su rol)
\item
  (Nombre del personaje y breve descripción)
\item
  (Nombre del personaje y breve descripción)
\end{itemize}

\textbf{Descripción del guion:}

\emph{(Escribir el desarrollo de la obra escena por escena, incluyendo
principales diálogos o narraciones. Cada escena debe reflejar ideas de
la teoría estudiada, con base a información real, \textbf{no
inventada})}

\textbf{Escena 1:} (Descripción breve + diálogos principales)

\textbf{Escena 2:} (Descripción breve + diálogos principales)

\textbf{Escena 3:} (Descripción breve + diálogos principales)

\textbf{Conclusión o mensaje final:} (Resumen del mensaje que la obra
deja al público sobre la teoría).

\textbf{Bibliografía:}

\textbf{Nota:}

\begin{itemize}
\item
  Cuidar ortografía y redacción.
\item
  Citar en las preguntas 1, 2 y 3.
\item
  El informe será evaluado junto a la presentación teatral sobre 2
  puntos.
\item
  La extensión recomendada del informe es de 3 a 5 páginas.
\end{itemize}

\textbf{Revisión final:} Antes de entregar, revisa que tu informe esté
ordenado, claro y completo.

\section{Mucho éxito en su investigación y
representación!}\label{mucho-uxe9xito-en-su-investigaciuxf3n-y-representaciuxf3n}




\end{document}
